% --- Archon physics formalization source begin ---
% archon:physics
% archon:covers PhyXMiniProblems/problem_phyx_mini_0723.lean
% archon:source-report reports/phyx_mini/problem_phyx_mini_0723.source.json

\chapter{Physics problem phyx\_mini\_0723}
\label{ch:PhyXMiniProblems_problem_phyx_mini_0723}

\paragraph{Problem source.}
\#\# Physical scenario

You need to determine the density of an unknown liquid. You notice that a block floats in this liquid with 4.6\textbackslash{} \textbackslash{}mathrm\{cm\} of the side of the block submerged. When the block is placed in water, it also floats but with 5.8\textbackslash{} \textbackslash{}mathrm\{cm\} submerged.

\#\# Question

What is the density of the unknown liquid?

\#\# Answer choices

- A: A: \$1360\textbackslash{}

- B: B: \$1160\textbackslash{}

- C: C: \$1260\textbackslash{}

- D: D: \$1200\textbackslash{}

\#\# Figure caption (auxiliary; use the image as primary evidence)

The image depicts two rectangular objects partially submerged in two different liquids. The objects and liquids are illustrated as follows:

1. **Left Side**: 
   - A rectangle is partially submerged in a liquid labeled "Unknown liquid."
   - The submerged depth of the rectangle is marked with a double-headed arrow labeled "h\_u."
   - The text "Submerged length" is written beside the arrow indicating depth.

2. **Right Side**:
   - Another rectangle, identical in shape to the first, is partially submerged in another liquid labeled "Water."
   - This rectangle’s submerged depth is marked with a double-headed arrow labeled "h\_w."

3. **Common Elements**:
   - Both rectangles share a label at the top denoting "Area A."
   - Horizontal lines represent the liquid surfaces, with one surface separating the unknown liquid and water.

The textured background differentiates the two liquids.

\#\# Dataset metadata

Statics; Physical Model Grounding Reasoning, Multi-Formula Reasoning

\paragraph{Recorded answer/context.}
C: C: \$1260\textbackslash{}

\paragraph{Figure/image path.}
/root/proposal\_for\_physic/science-mango/phyx\_mini\_run/phyx\_data/test\_image/723.png

\paragraph{Formalization target.}
create a compiling Lean file with sorry bodies at `PhyXMiniProblems/problem\_phyx\_mini\_0723.lean`. The Lean declarations must preserve the physical quantities, dimensions or dimensional roles, figure labels, governing-law hypotheses, and final relation expressed by this problem.
Use Mathlib/Physlib names found through LeanExplore where available. If a physics API is missing, introduce faithful local abstractions rather than scalar placeholder aliases.

\begin{theorem}[Physics formalization target]
\label{thm:physics:phyx_mini_0723:target}
\lean{PhyXMiniProblems.ProblemPhyXMini0723.problem_phyx_mini_0723}
\uses{def:physics:phyx-mini-0723:phyxminiproblems-problemphyxmini0723-densityinkilogramspercubicmeter, def:physics:phyx-mini-0723:phyxminiproblems-problemphyxmini0723-floatingblocksetup, def:physics:phyx-mini-0723:phyxminiproblems-problemphyxmini0723-matchesproblemandfiguredata, def:physics:phyx-mini-0723:phyxminiproblems-problemphyxmini0723-usesfreshwaterdensity, def:physics:phyx-mini-0723:phyxminiproblems-problemphyxmini0723-hasphysicalparameters, def:physics:phyx-mini-0723:phyxminiproblems-problemphyxmini0723-satisfiesprismaticdisplacementlaw, def:physics:phyx-mini-0723:phyxminiproblems-problemphyxmini0723-satisfiesarchimedesprinciple, def:physics:phyx-mini-0723:phyxminiproblems-problemphyxmini0723-floatsinstaticequilibrium, def:physics:phyx-mini-0723:phyxminiproblems-problemphyxmini0723-answerchoice, def:physics:phyx-mini-0723:phyxminiproblems-problemphyxmini0723-isclosestdisplayeddensity}
The assigned autoformalize agent should translate this physics problem into Lean declarations in the covered file, with theorem and lemma proof bodies written as `by sorry`.
\end{theorem}
\begin{proof}
This is an autoformalization task, not a proof task. Produce statements that can later be proved by the physics prover without weakening the physical model.
\end{proof}

% --- Archon declaration topology begin ---
\section{Declaration topology}
\begin{definition}[Length Quantity]
  \label{def:physics:phyx-mini-0723:phyxminiproblems-problemphyxmini0723-lengthquantity}
  \lean{PhyXMiniProblems.ProblemPhyXMini0723.LengthQuantity}
  A nonnegative physical length
\end{definition}
\begin{proof}
  This is the declared model definition of Length Quantity.
\end{proof}
\begin{definition}[Area Quantity]
  \label{def:physics:phyx-mini-0723:phyxminiproblems-problemphyxmini0723-areaquantity}
  \lean{PhyXMiniProblems.ProblemPhyXMini0723.AreaQuantity}
  A nonnegative physical area, with dimension L²
\end{definition}
\begin{proof}
  This is the declared model definition of Area Quantity.
\end{proof}
\begin{definition}[Volume Quantity]
  \label{def:physics:phyx-mini-0723:phyxminiproblems-problemphyxmini0723-volumequantity}
  \lean{PhyXMiniProblems.ProblemPhyXMini0723.VolumeQuantity}
  A nonnegative physical volume, with dimension L³
\end{definition}
\begin{proof}
  This is the declared model definition of Volume Quantity.
\end{proof}
\begin{definition}[Mass Density Quantity]
  \label{def:physics:phyx-mini-0723:phyxminiproblems-problemphyxmini0723-massdensityquantity}
  \lean{PhyXMiniProblems.ProblemPhyXMini0723.MassDensityQuantity}
  A nonnegative mass density, with dimension M L⁻³
\end{definition}
\begin{proof}
  This is the declared model definition of Mass Density Quantity.
\end{proof}
\begin{definition}[Acceleration Quantity]
  \label{def:physics:phyx-mini-0723:phyxminiproblems-problemphyxmini0723-accelerationquantity}
  \lean{PhyXMiniProblems.ProblemPhyXMini0723.AccelerationQuantity}
  A nonnegative acceleration magnitude, with dimension L T⁻²
\end{definition}
\begin{proof}
  This is the declared model definition of Acceleration Quantity.
\end{proof}
\begin{definition}[Force Quantity]
  \label{def:physics:phyx-mini-0723:phyxminiproblems-problemphyxmini0723-forcequantity}
  \lean{PhyXMiniProblems.ProblemPhyXMini0723.ForceQuantity}
  A nonnegative force magnitude, with dimension M L T⁻²
\end{definition}
\begin{proof}
  This is the declared model definition of Force Quantity.
\end{proof}
\begin{definition}[nonnegative SIReadout]
  \label{def:physics:phyx-mini-0723:phyxminiproblems-problemphyxmini0723-nonnegativesireadout}
  \lean{PhyXMiniProblems.ProblemPhyXMini0723.nonnegativeSIReadout}
  Read a nonnegative dimensionful quantity in coherent SI units
\end{definition}
\begin{proof}
  This is the declared model definition of nonnegative SIReadout.
\end{proof}
\begin{definition}[length In Meters]
  \label{def:physics:phyx-mini-0723:phyxminiproblems-problemphyxmini0723-lengthinmeters}
  \lean{PhyXMiniProblems.ProblemPhyXMini0723.lengthInMeters}
  \uses{def:physics:phyx-mini-0723:phyxminiproblems-problemphyxmini0723-lengthquantity, def:physics:phyx-mini-0723:phyxminiproblems-problemphyxmini0723-nonnegativesireadout}
  Numerical readout of a physical length in meters
\end{definition}
\begin{proof}
  This is the declared model definition of length In Meters.
\end{proof}
\begin{definition}[length In Centimeters]
  \label{def:physics:phyx-mini-0723:phyxminiproblems-problemphyxmini0723-lengthincentimeters}
  \lean{PhyXMiniProblems.ProblemPhyXMini0723.lengthInCentimeters}
  \uses{def:physics:phyx-mini-0723:phyxminiproblems-problemphyxmini0723-lengthquantity}
  Numerical readout of a physical length in centimeters
\end{definition}
\begin{proof}
  This is the declared model definition of length In Centimeters.
\end{proof}
\begin{definition}[area In Square Meters]
  \label{def:physics:phyx-mini-0723:phyxminiproblems-problemphyxmini0723-areainsquaremeters}
  \lean{PhyXMiniProblems.ProblemPhyXMini0723.areaInSquareMeters}
  \uses{def:physics:phyx-mini-0723:phyxminiproblems-problemphyxmini0723-areaquantity, def:physics:phyx-mini-0723:phyxminiproblems-problemphyxmini0723-nonnegativesireadout}
  Numerical readout of a physical area in square meters
\end{definition}
\begin{proof}
  This is the declared model definition of area In Square Meters.
\end{proof}
\begin{definition}[volume In Cubic Meters]
  \label{def:physics:phyx-mini-0723:phyxminiproblems-problemphyxmini0723-volumeincubicmeters}
  \lean{PhyXMiniProblems.ProblemPhyXMini0723.volumeInCubicMeters}
  \uses{def:physics:phyx-mini-0723:phyxminiproblems-problemphyxmini0723-volumequantity, def:physics:phyx-mini-0723:phyxminiproblems-problemphyxmini0723-nonnegativesireadout}
  Numerical readout of a physical volume in cubic meters
\end{definition}
\begin{proof}
  This is the declared model definition of volume In Cubic Meters.
\end{proof}
\begin{definition}[density In Kilograms Per Cubic Meter]
  \label{def:physics:phyx-mini-0723:phyxminiproblems-problemphyxmini0723-densityinkilogramspercubicmeter}
  \lean{PhyXMiniProblems.ProblemPhyXMini0723.densityInKilogramsPerCubicMeter}
  \uses{def:physics:phyx-mini-0723:phyxminiproblems-problemphyxmini0723-massdensityquantity, def:physics:phyx-mini-0723:phyxminiproblems-problemphyxmini0723-nonnegativesireadout}
  Numerical readout of a mass density in kilograms per cubic meter
\end{definition}
\begin{proof}
  This is the declared model definition of density In Kilograms Per Cubic Meter.
\end{proof}
\begin{definition}[acceleration In Meters Per Second Squared]
  \label{def:physics:phyx-mini-0723:phyxminiproblems-problemphyxmini0723-accelerationinmeterspersecondsquared}
  \lean{PhyXMiniProblems.ProblemPhyXMini0723.accelerationInMetersPerSecondSquared}
  \uses{def:physics:phyx-mini-0723:phyxminiproblems-problemphyxmini0723-accelerationquantity, def:physics:phyx-mini-0723:phyxminiproblems-problemphyxmini0723-nonnegativesireadout}
  Numerical readout of an acceleration in meters per second squared
\end{definition}
\begin{proof}
  This is the declared model definition of acceleration In Meters Per Second Squared.
\end{proof}
\begin{definition}[force In Newtons]
  \label{def:physics:phyx-mini-0723:phyxminiproblems-problemphyxmini0723-forceinnewtons}
  \lean{PhyXMiniProblems.ProblemPhyXMini0723.forceInNewtons}
  \uses{def:physics:phyx-mini-0723:phyxminiproblems-problemphyxmini0723-forcequantity, def:physics:phyx-mini-0723:phyxminiproblems-problemphyxmini0723-nonnegativesireadout}
  Numerical readout of a force in newtons
\end{definition}
\begin{proof}
  This is the declared model definition of force In Newtons.
\end{proof}
\begin{definition}[Placement]
  \label{def:physics:phyx-mini-0723:phyxminiproblems-problemphyxmini0723-placement}
  \lean{PhyXMiniProblems.ProblemPhyXMini0723.Placement}
  The two placements of the same block shown in the figure
\end{definition}
\begin{proof}
  This is the declared model definition of Placement.
\end{proof}
\begin{definition}[Figure Liquid Label]
  \label{def:physics:phyx-mini-0723:phyxminiproblems-problemphyxmini0723-figureliquidlabel}
  \lean{PhyXMiniProblems.ProblemPhyXMini0723.FigureLiquidLabel}
  Liquid names printed on the two panels of the primary figure
\end{definition}
\begin{proof}
  This is the declared model definition of Figure Liquid Label.
\end{proof}
\begin{definition}[Figure Depth Symbol]
  \label{def:physics:phyx-mini-0723:phyxminiproblems-problemphyxmini0723-figuredepthsymbol}
  \lean{PhyXMiniProblems.ProblemPhyXMini0723.FigureDepthSymbol}
  Symbols attached to the submerged-length arrows in the figure
\end{definition}
\begin{proof}
  This is the declared model definition of Figure Depth Symbol.
\end{proof}
\begin{definition}[Figure Area Symbol]
  \label{def:physics:phyx-mini-0723:phyxminiproblems-problemphyxmini0723-figureareasymbol}
  \lean{PhyXMiniProblems.ProblemPhyXMini0723.FigureAreaSymbol}
  The common top-face-area symbol printed over both block drawings
\end{definition}
\begin{proof}
  This is the declared model definition of Figure Area Symbol.
\end{proof}
\begin{definition}[Supplied Figure Readout]
  \label{def:physics:phyx-mini-0723:phyxminiproblems-problemphyxmini0723-suppliedfigurereadout}
  \lean{PhyXMiniProblems.ProblemPhyXMini0723.SuppliedFigureReadout}
  \uses{def:physics:phyx-mini-0723:phyxminiproblems-problemphyxmini0723-lengthquantity, def:physics:phyx-mini-0723:phyxminiproblems-problemphyxmini0723-areaquantity, def:physics:phyx-mini-0723:phyxminiproblems-problemphyxmini0723-placement, def:physics:phyx-mini-0723:phyxminiproblems-problemphyxmini0723-figureliquidlabel, def:physics:phyx-mini-0723:phyxminiproblems-problemphyxmini0723-figuredepthsymbol, def:physics:phyx-mini-0723:phyxminiproblems-problemphyxmini0723-figureareasymbol}
  Quantitative and categorical information read from the supplied image
\end{definition}
\begin{proof}
  This is the declared model definition of Supplied Figure Readout.
\end{proof}
\begin{definition}[Floating Block Setup]
  \label{def:physics:phyx-mini-0723:phyxminiproblems-problemphyxmini0723-floatingblocksetup}
  \lean{PhyXMiniProblems.ProblemPhyXMini0723.FloatingBlockSetup}
  \uses{def:physics:phyx-mini-0723:phyxminiproblems-problemphyxmini0723-lengthquantity, def:physics:phyx-mini-0723:phyxminiproblems-problemphyxmini0723-areaquantity, def:physics:phyx-mini-0723:phyxminiproblems-problemphyxmini0723-volumequantity, def:physics:phyx-mini-0723:phyxminiproblems-problemphyxmini0723-massdensityquantity, def:physics:phyx-mini-0723:phyxminiproblems-problemphyxmini0723-accelerationquantity, def:physics:phyx-mini-0723:phyxminiproblems-problemphyxmini0723-forcequantity, def:physics:phyx-mini-0723:phyxminiproblems-problemphyxmini0723-placement, def:physics:phyx-mini-0723:phyxminiproblems-problemphyxmini0723-suppliedfigurereadout}
  The physical state used for both placements of the same block. The single fields crossSectionalArea, gravitationalAcceleration, and blockWeight encode that the same prismatic block and the same gravitational environment are used in both liquids. The requested unknown density remains an independent component of fluidDensity; it is not defined from the answer
\end{definition}
\begin{proof}
  This is the declared model definition of Floating Block Setup.
\end{proof}
\begin{definition}[Matches Problem And Figure Data]
  \label{def:physics:phyx-mini-0723:phyxminiproblems-problemphyxmini0723-matchesproblemandfiguredata}
  \lean{PhyXMiniProblems.ProblemPhyXMini0723.MatchesProblemAndFigureData}
  \uses{def:physics:phyx-mini-0723:phyxminiproblems-problemphyxmini0723-lengthincentimeters, def:physics:phyx-mini-0723:phyxminiproblems-problemphyxmini0723-floatingblocksetup}
  Data stated in the prose or read from the primary figure. The decimal depth readouts are written as exact rational centimeter values
\end{definition}
\begin{proof}
  This is the declared model definition of Matches Problem And Figure Data.
\end{proof}
\begin{definition}[Uses Fresh Water Density]
  \label{def:physics:phyx-mini-0723:phyxminiproblems-problemphyxmini0723-usesfreshwaterdensity}
  \lean{PhyXMiniProblems.ProblemPhyXMini0723.UsesFreshWaterDensity}
  \uses{def:physics:phyx-mini-0723:phyxminiproblems-problemphyxmini0723-densityinkilogramspercubicmeter, def:physics:phyx-mini-0723:phyxminiproblems-problemphyxmini0723-floatingblocksetup}
  The conventional fresh-water density used by the recorded answer
\end{definition}
\begin{proof}
  This is the declared model definition of Uses Fresh Water Density.
\end{proof}
\begin{definition}[Has Physical Parameters]
  \label{def:physics:phyx-mini-0723:phyxminiproblems-problemphyxmini0723-hasphysicalparameters}
  \lean{PhyXMiniProblems.ProblemPhyXMini0723.HasPhysicalParameters}
  \uses{def:physics:phyx-mini-0723:phyxminiproblems-problemphyxmini0723-areainsquaremeters, def:physics:phyx-mini-0723:phyxminiproblems-problemphyxmini0723-accelerationinmeterspersecondsquared, def:physics:phyx-mini-0723:phyxminiproblems-problemphyxmini0723-floatingblocksetup}
  Nondegeneracy assumptions needed to compare the two force balances
\end{definition}
\begin{proof}
  This is the declared model definition of Has Physical Parameters.
\end{proof}
\begin{definition}[Satisfies Prismatic Displacement Law]
  \label{def:physics:phyx-mini-0723:phyxminiproblems-problemphyxmini0723-satisfiesprismaticdisplacementlaw}
  \lean{PhyXMiniProblems.ProblemPhyXMini0723.SatisfiesPrismaticDisplacementLaw}
  \uses{def:physics:phyx-mini-0723:phyxminiproblems-problemphyxmini0723-lengthinmeters, def:physics:phyx-mini-0723:phyxminiproblems-problemphyxmini0723-areainsquaremeters, def:physics:phyx-mini-0723:phyxminiproblems-problemphyxmini0723-volumeincubicmeters, def:physics:phyx-mini-0723:phyxminiproblems-problemphyxmini0723-placement, def:physics:phyx-mini-0723:phyxminiproblems-problemphyxmini0723-floatingblocksetup}
  For the constant-cross-section block, displaced volume is top-face area times submerged depth in either placement
\end{definition}
\begin{proof}
  This is the declared model definition of Satisfies Prismatic Displacement Law.
\end{proof}
\begin{definition}[Satisfies Archimedes Principle]
  \label{def:physics:phyx-mini-0723:phyxminiproblems-problemphyxmini0723-satisfiesarchimedesprinciple}
  \lean{PhyXMiniProblems.ProblemPhyXMini0723.SatisfiesArchimedesPrinciple}
  \uses{def:physics:phyx-mini-0723:phyxminiproblems-problemphyxmini0723-volumeincubicmeters, def:physics:phyx-mini-0723:phyxminiproblems-problemphyxmini0723-densityinkilogramspercubicmeter, def:physics:phyx-mini-0723:phyxminiproblems-problemphyxmini0723-accelerationinmeterspersecondsquared, def:physics:phyx-mini-0723:phyxminiproblems-problemphyxmini0723-forceinnewtons, def:physics:phyx-mini-0723:phyxminiproblems-problemphyxmini0723-placement, def:physics:phyx-mini-0723:phyxminiproblems-problemphyxmini0723-floatingblocksetup}
  Archimedes' law: buoyancy is fluid density × displaced volume × g
\end{definition}
\begin{proof}
  This is the declared model definition of Satisfies Archimedes Principle.
\end{proof}
\begin{definition}[Floats In Static Equilibrium]
  \label{def:physics:phyx-mini-0723:phyxminiproblems-problemphyxmini0723-floatsinstaticequilibrium}
  \lean{PhyXMiniProblems.ProblemPhyXMini0723.FloatsInStaticEquilibrium}
  \uses{def:physics:phyx-mini-0723:phyxminiproblems-problemphyxmini0723-forceinnewtons, def:physics:phyx-mini-0723:phyxminiproblems-problemphyxmini0723-placement, def:physics:phyx-mini-0723:phyxminiproblems-problemphyxmini0723-floatingblocksetup}
  Static flotation: in each liquid, buoyancy balances the same block weight
\end{definition}
\begin{proof}
  This is the declared model definition of Floats In Static Equilibrium.
\end{proof}
\begin{definition}[Answer Choice]
  \label{def:physics:phyx-mini-0723:phyxminiproblems-problemphyxmini0723-answerchoice}
  \lean{PhyXMiniProblems.ProblemPhyXMini0723.AnswerChoice}
  Labels of the four density choices printed with the problem
\end{definition}
\begin{proof}
  This is the declared model definition of Answer Choice.
\end{proof}
\begin{definition}[density In Kilograms Per Cubic Meter]
  \label{def:physics:phyx-mini-0723:phyxminiproblems-problemphyxmini0723-answerchoice-densityinkilogramspercubicmeter}
  \lean{PhyXMiniProblems.ProblemPhyXMini0723.AnswerChoice.densityInKilogramsPerCubicMeter}
  \uses{def:physics:phyx-mini-0723:phyxminiproblems-problemphyxmini0723-densityinkilogramspercubicmeter, def:physics:phyx-mini-0723:phyxminiproblems-problemphyxmini0723-answerchoice}
  Density readout in kilograms per cubic meter printed for each choice
\end{definition}
\begin{proof}
  This is the declared model definition of density In Kilograms Per Cubic Meter.
\end{proof}
\begin{definition}[recorded Answer Choice]
  \label{def:physics:phyx-mini-0723:phyxminiproblems-problemphyxmini0723-recordedanswerchoice}
  \lean{PhyXMiniProblems.ProblemPhyXMini0723.recordedAnswerChoice}
  \uses{def:physics:phyx-mini-0723:phyxminiproblems-problemphyxmini0723-answerchoice}
  The answer key supplied with the dataset, retained only as metadata
\end{definition}
\begin{proof}
  This is the declared model definition of recorded Answer Choice.
\end{proof}
\begin{definition}[Is Closest Displayed Density]
  \label{def:physics:phyx-mini-0723:phyxminiproblems-problemphyxmini0723-isclosestdisplayeddensity}
  \lean{PhyXMiniProblems.ProblemPhyXMini0723.IsClosestDisplayedDensity}
  \uses{def:physics:phyx-mini-0723:phyxminiproblems-problemphyxmini0723-densityinkilogramspercubicmeter, def:physics:phyx-mini-0723:phyxminiproblems-problemphyxmini0723-floatingblocksetup, def:physics:phyx-mini-0723:phyxminiproblems-problemphyxmini0723-answerchoice}
  A displayed density choice is at least as close as every alternative
\end{definition}
\begin{proof}
  This is the declared model definition of Is Closest Displayed Density.
\end{proof}
\begin{lemma}[density times depth eq water density times depth]
  \label{lem:physics:phyx-mini-0723:phyxminiproblems-problemphyxmini0723-density-times-depth-eq-water-density-times-depth}
  \lean{PhyXMiniProblems.ProblemPhyXMini0723.density_times_depth_eq_water_density_times_depth}
  \uses{def:physics:phyx-mini-0723:phyxminiproblems-problemphyxmini0723-lengthinmeters, def:physics:phyx-mini-0723:phyxminiproblems-problemphyxmini0723-densityinkilogramspercubicmeter, def:physics:phyx-mini-0723:phyxminiproblems-problemphyxmini0723-floatingblocksetup, def:physics:phyx-mini-0723:phyxminiproblems-problemphyxmini0723-hasphysicalparameters, def:physics:phyx-mini-0723:phyxminiproblems-problemphyxmini0723-satisfiesprismaticdisplacementlaw, def:physics:phyx-mini-0723:phyxminiproblems-problemphyxmini0723-satisfiesarchimedesprinciple, def:physics:phyx-mini-0723:phyxminiproblems-problemphyxmini0723-floatsinstaticequilibrium}
  Comparing the two equilibria cancels the common area, gravity, and block weight, leaving the standard density-depth relation
\end{lemma}
\begin{proof}
  The conclusion follows from the governing relations and figure readouts named in the statement of density times depth eq water density times depth.
\end{proof}
\begin{lemma}[unknown liquid density exact]
  \label{lem:physics:phyx-mini-0723:phyxminiproblems-problemphyxmini0723-unknown-liquid-density-exact}
  \lean{PhyXMiniProblems.ProblemPhyXMini0723.unknown_liquid_density_exact}
  \uses{def:physics:phyx-mini-0723:phyxminiproblems-problemphyxmini0723-densityinkilogramspercubicmeter, def:physics:phyx-mini-0723:phyxminiproblems-problemphyxmini0723-floatingblocksetup, def:physics:phyx-mini-0723:phyxminiproblems-problemphyxmini0723-matchesproblemandfiguredata, def:physics:phyx-mini-0723:phyxminiproblems-problemphyxmini0723-usesfreshwaterdensity, def:physics:phyx-mini-0723:phyxminiproblems-problemphyxmini0723-hasphysicalparameters, def:physics:phyx-mini-0723:phyxminiproblems-problemphyxmini0723-satisfiesprismaticdisplacementlaw, def:physics:phyx-mini-0723:phyxminiproblems-problemphyxmini0723-satisfiesarchimedesprinciple, def:physics:phyx-mini-0723:phyxminiproblems-problemphyxmini0723-floatsinstaticequilibrium}
  The exact ideal-model density is 1000 * 5.8 / 4.6 = 29000 / 23 kg/m³
\end{lemma}
\begin{proof}
  The conclusion follows from the governing relations and figure readouts named in the statement of unknown liquid density exact.
\end{proof}
% --- Archon declaration topology end ---
% --- Archon physics formalization source end ---
